\documentclass[11pt, reqno]{amsart}

%%% Font Setup
% \usepackage{ETbb}
\usepackage{libertinus}           % text + math, type1 so pdflatex-safe
\usepackage[T1]{fontenc}

%%% Packages
\usepackage{amsmath, amssymb, amsthm}
\usepackage{xcolor}
\usepackage{bm}
\usepackage[colorlinks=true, linkcolor=blue!60!black, citecolor=blue!60!black, urlcolor=blue!60!black]{hyperref}
\usepackage{enumitem}
\usepackage{graphicx}
\usepackage{subcaption}
\usepackage{flafter}

\usepackage[framemethod=TikZ]{mdframed}
\mdfdefinestyle{lecturebox}{
  backgroundcolor=black!5,      % light gray fill
  linecolor=black!12,           % faint border; set linewidth=0pt to drop it
  linewidth=0.6pt,
  roundcorner=3pt,
  leftmargin=1em, rightmargin=1em,        % outer indent from text block
  innerleftmargin=1.4em, innerrightmargin=1.4em,
  innertopmargin=1em, innerbottommargin=1em,
}

%%% Page Geometry
\usepackage{geometry}
\geometry{margin=1.2in, top=1in, bottom=1in}

%%% Section Headings in Small Caps
\makeatletter
\def\section{\@startsection{section}{1}%
  \z@{.7\linespacing\@plus\linespacing}{.5\linespacing}%
  {\normalfont\scshape\centering}}
\makeatother

%%% Title in Small Caps
\makeatletter
\def\@settitle{\begin{center}%
  \baselineskip14\p@\relax
  \bfseries\scshape\uppercasenonmath\@title
  \@title
  \end{center}%
}
\makeatother

%%% Theorem Environments
\newtheoremstyle{notestyle}
  {6pt}{6pt}{\itshape}{}{\bfseries}{.}{.5em}
  {\thmname{#1}\thmnumber{ #2}\thmnote{ \normalfont(#3)}}

\theoremstyle{notestyle}
\newtheorem{theorem}{Theorem}[section]
\newtheorem{lemma}[theorem]{Lemma}
\newtheorem{proposition}[theorem]{Proposition}
\newtheorem{corollary}[theorem]{Corollary}

\newtheoremstyle{defstyle}
  {6pt}{6pt}{\normalfont}{}{\bfseries}{.}{.5em}
  {\thmname{#1}\thmnumber{ #2}\thmnote{ \normalfont(#3)}}

\theoremstyle{defstyle}
\newtheorem{definition}[theorem]{Definition}
\newtheorem{example}[theorem]{Example}
\newtheorem{remark}[theorem]{Remark}

\newcommand{\R}{\mathbb{R}}
\newcommand{\T}{^\top}
\newcommand{\mc}[1]{\mathcal{#1}}

%%% Additional notation
\newcommand{\E}{\mathbb{E}}
\newcommand{\V}{\mathcal{V}}
\newcommand{\Hphys}{\mathcal{H}_{\mathrm{phys}}}
\newcommand{\Cov}{\operatorname{Cov}}
\newcommand{\Tr}{\operatorname{Tr}}
\newcommand{\Id}{\mathrm{Id}}
\newcommand{\argmin}{\operatorname*{arg\,min}}
\newcommand{\norm}[1]{\left\lVert #1 \right\rVert}
\newcommand{\inner}[2]{\left\langle #1,#2 \right\rangle}
\newcommand{\MPS}{\mathcal{M}_{\chi}}
\newcommand{\Q}{\mathcal{Q}}
\newcommand{\FS}{\mathrm{FS}}
\newcommand{\Heff}{H_{\mathrm{eff}}}
\newcommand{\Neff}{N_{\mathrm{eff}}}

%%% Metadata
\title{Correlated Cubic Features}
\author{authors}
\date{}

\begin{document}
\maketitle
\thispagestyle{empty}

\noindent\textit{%
We need to extend Nisch's two-feature calculation to $k$ correlated target features at Bayesian equilibrium. In particular, when does one committee serve all the features, when does it split, and what populations of neurons minimize the posterior cost after splitting?}

% \tableofcontents

% \newpage

\section{Setup}

Let $N$ be the input dimension, $K$ be the number of hidden neurons, and $P$ be the number of training examples.

Take
\begin{equation}
    K = \frac{N^2}{\Delta},\quad P = \chi K,
\end{equation}
with $\Delta, \chi > 0$ fixed. Our sample ratio remains
\begin{equation}
    \alpha = \frac PN = \frac{\chi N}{\Delta},
\end{equation}
growing with $N$. The source's sample ratio is $\chi = P/K$, not our $\alpha$.

\section{Teacher and Training Data}
Draw independent $\mathbf x^\mu \sim \mathcal N(0, I_N)$. Choose unit teacher directions satisfying
\begin{equation}
    \mathbf e_\ell^\top \mathbf e_t = G_{\ell t} = (1 - \omega)d_{\ell t} + \omega,\quad 0 \leq \omega \leq 1.
\end{equation}
Initially use noiseless labels
\begin{equation}
    y^\mu = f_*(\mathbf x^\mu),\quad f_*(\mathbf x) = B\sum_{\ell = 1}^k \psi_3(\mathbf e_\ell^\top \mathbf x),\quad \psi_3(z) = \frac{H_3(z)}{\sqrt 6} = \frac{z^3 - 3z}{\sqrt{6}}.
\end{equation}
Where $H_3$ is the third Hermite polynomial. The normalization makes each feature have unit variance. For Gaussian inputs,
\begin{equation}
    \mathbb E_{\mathbf x} [\psi_3(\mathbf e_\ell^\top \mathbf x)\psi_3(\mathbf e_t^\top\mathbf x)] = G_{\ell t}^3.
\end{equation}
Thus $\omega$ is a direction overlap. The corresponding cubic-feature correlation is $\omega^3$.
\section{Student network and equilibrium ensemble}
Use the pure cubic student
\begin{equation}
    f_\theta(\mathbf x) = \sum_{i=1}^K a_i \gamma_3 \psi_3\left( \frac{\mathbf w_i^\top \mathbf x}{\sqrt N}\right),\quad \gamma_3 > 0.
\end{equation}
For this calculation, the prior is uniform over hidden-weight spheres and over readouts satisfying:
\begin{equation}
    \|\mathbf w_i\|^2 = N,\quad |a_i| \leq \frac{A}{\sqrt K},\quad \sum_{i=1}^K a_i^2 = \tau^2,\quad 0 < \tau^2 < A^2.
\end{equation}
Here, $A, \tau, \gamma_3$ are fixed. The factor $K^{-1/2}$ in Nisch's network has been absorbed into our readouts $a_i$. For fixed dataset $\mathcal D$, define
\begin{equation}
    E_{\mathcal D}(\theta) = \frac 12 \sum_{\mu=1}^P[y^\mu - f_\theta(\mathbf x^\mu)]^2,\quad p_\beta(\theta \mid \mathcal D) = \frac{p_0(\theta)e^{-\beta E_{\mathcal D}(\theta)}}{Z_\beta (\mathcal D)}.
\end{equation}
This is an equilibrium problem. It doesn't specify how quickly SGD or a posterior sampler reaches the predicted populations. We write $\langle \cdot \rangle_{\mathcal D}$ for posterior averages and $\mathbb E_{\mathbf x}$ for averages over fresh Gaussian inputs.

Define $\mathbf w_\ell^* = \sqrt N \mathbf e_\ell$. The signed teacher overlaps are
\begin{equation}
    m_{i\ell} = \frac{\mathbf w_i^\top \mathbf w_\ell^*}{N},\quad q_i = \frac{\|\mathbf w_i\|^2}{N} = 1.
\end{equation}
\section{Reduced description in terms of committees}
Let $\mathcal V = \mathrm{span}\{\mathbf e_1, \ldots, \mathbf e_k\}$. A committee consists of neurons sharing a projection into $\mathcal V$:
\begin{equation}
    \frac{\mathbf w_i}{\sqrt N} = r_c \mathbf u_c + \sqrt{1 - r_c^2 }\bm \xi_i,\quad \mathbf u_c \in \mathcal V,\quad \|\mathbf u_c\| = \|\bm \xi_i\| = 1,\quad \bm\xi_i \perp \mathcal V.
\end{equation}
$\mathcal V$ is a small subspace inside the input space such that the teacher depends only on the projections of $\mathbf x$ into $\mathcal V$. We can therefore split a neuron's normalized weight into the above decomposition. The first term is the projection into the teacher subspace, the second term is the transverse component. Here, $\mathbf u_c$ tells us which direction the group represents, while $r_c$ tells us how \textit{strongly neurons concentrates on that directions.} A committee shares the first component; its neurons can differ in their transverse components. The teacher overlaps are $m_{i\ell} = r_c \mathbf u_c^\top \mathbf e_\ell$.

The useful identity here is
\begin{equation}
    \mathbb E_Z \left[ \psi_3 (rz + \sqrt{1-r^2}Z)\right] = r^3 \psi_3(z).
\end{equation}
So part of the neuron's response that survives averaging over transverse inputs is a cubic feature along $\mathbf u_c$, with strength proportional to $r_c^3$. 

Under the spherical prior, most weight distributions are almost orthogonal to any fixed low-dim subspace. A finite projection $r$ is rare because comparitively few directions have that projection:
\begin{equation}
    \mathrm{Pr}(r) \asymp e^{-N J(r)},\quad J(r) = -\frac 12\log(1 - r^2).
\end{equation}
This is an \textit{entropy cost}. Every allowed direction is equally likely, but there are vastly more directions with small project than large.

Stronger alignment gives more cubic output per neuron (prop to $r^3$), but costs more prior probability. 

Nisch' leading-order calculation minimizes the alignment cost per unit cubic output:
\begin{equation}
    \frac{\text{alignment cost}}{\text{cubic output}} \propto \frac{J(r)}{r^3},\quad r_* = \argmin_{0 < r < 1} \frac{J(r)}{r^3} \simeq 0.73047.
\end{equation}
So weak alignment produces too little signal; alignment approaching one becomes too statistically costly. 

After optimizing readout magnitudes and projection lengths, represent the retained cubic signal as
\begin{equation}
    f_{\mathcal C}(\mathbf x) = \sum_{c=1}^L b_c \psi_3(\mathbf u_c^\top \mathbf x),\quad b_c \geq 0.
\end{equation}
\textit{The number $L$ of committees is unknown.} Signs can also be absorbed into $\mathbf u_c$, since $\psi_3$ is odd.

Under the background approximation, the effective population cost is
\begin{equation}
    F(\mathcal C) = \lambda \sum_c b_c + \frac 12 \mathbb E_{\mathbf x}\left[ f_*(\mathbf x) - f_{\mathcal C}(\mathbf x)\right]^2
\end{equation}
with 
\begin{equation}
    \lambda = \frac{\sqrt{\Delta}J(r_*)}{g_{\text{eff}}\gamma_3 A r_*^3}.
\end{equation}

\section{Three Correlated Features}

Similarly to the two correlated features notes, let
\begin{align*}
f_* (x) = \beta \sum_{a=1}^3 \psi_3(e_a^\top x),
\qquad e_a^\top e_b = \omega \quad (a \neq b)
\end{align*}
Thus, we have three features, with mutual inner products all of $\omega$.
This case will be easier than the four-feature asymmetric case; not only do we have one fewer feature, there is also an $S_3$ symmetry.

We can first rederive the free energy expression here, in a very similar way to the two feature case.
If $w$ is one hidden weight, we denote by $p$ its projection into $S = e_1 \times e_2 \times e_3$ space, and $u$ the unit direction of that projection:
\begin{align*}
w = pu + \sqrt{1-p^2} w_\perp
\end{align*}
As in the two-feature calculation,
\begin{align*}
\mathbb{E}_{x_\perp} \psi_3(w^\top x)
= p^3 \psi_3 (u^\top x).
\end{align*}
We again have that we need to minimize the alignment cost, which requires minimizing
\begin{align*}
\frac{J(p)}{p},
\qquad J(p) = -\frac{1}{2} \log (1-p^2)
\end{align*}
This all results in the same preferred radius $p = r_3 \approx 0.73$, and the same optimal $\lambda$.
Thus, these values do not depend on the number of features.

The ``free energy" that ends up being minimized is written
\begin{align*}
F = \lambda \sum_{j} b_j + \frac{1}{2} \mathbb{E}_x [f_*(x) - f_C(x)]^2
\qquad f_C(x) = \sum_j b_j \psi_3(u_j^\top x)
\end{align*}
as before.

Now, we perform a similar centering step as in the two-feature case, this time around a central axis, in light of the $S_3$ symmetry.
Define the normalized symmetric axis vector
\begin{align*}
e_0 = \frac{e_1 + e_2 + e_3}{\sqrt{3(1+2\omega)}}
\end{align*}
which lets us write any of the teacher vectors as
\begin{align*}
    e_a = c e_0 + s v_a,
    \qquad c = \sqrt{\frac{1+2\omega}{3}},
    \qquad s = \sqrt{\frac{2(1-\omega)}{3}}
\end{align*}
where the $v_a$ are three unit vectors in the plane perpendicular to $e_0$, each pointing towards $e_a$, which form an equilateral triangle in that plane.

We can parametrize the projection of $w$, a committee representative, as
\begin{align*}
    w = pu + \sqrt{1-p^2}w_\perp,
    \qquad u = t e_0 + \sqrt{1-t^2} v
\end{align*}
with $v$ lying in the $e_0$ perpendicular plane.

Now, similarly to the two feature case, we look at the difference in free energy from adding a single committee with output $b \psi_3 (u^\top x)$ to a completely unaligned network with $f_C = 0$:
\begin{align*}
    F(b, u) - F(0)
    &= b \left[\lambda - \mathbb{E}_x \{f_* \psi_3 (u^\top x)\} \right]
    + \frac{b^2}{2} \mathbb{E}_x \psi_3 (u^\top x)^2 \\
    &= b \left[ \lambda - \beta \{ (e_1^\top x)^3 + (e_2^\top x)^3 + (e_3^\top x)^3 \} 
    \right] + \frac{b^2}{2}
\end{align*}
where we used Hermite orthogonality to evaluate the expectations.
Then define
\begin{align*}
    H(t, \theta) = \sum_{a=1}^3 (e_a^\top u)^3
\end{align*}
where we have written on the left side the dependence of $H$ on $t$ and $\theta$, which we define as the angle between $v$ and $v_1$.
After some algebraic magic happens, it can be shown that
\begin{align}\label{eq:three-feature-H}
    H(t, \theta)
    &= 3c^3 t^3
    + \frac{9}{2} cs^2 t (1-t^2)
    + \frac{3}{4} s^3 (1-t^2)^{3/2} \cos (3\theta)
\end{align}
Recall that $e_a = c e_0 + s v_a$ and $u = t e_0 + \sqrt{1-t^2}v$.
Schematically, then, each of these terms in $H$ comes from the different combinations of terms in $e_a$ and $u$.
The first corresponds to three $e_0$ components, the second to one $e_0$ and two transverse components.
It turns out the term from two $e_0$ and one transverse vanishes because $v_1 + v_2 + v_3 = 0$. 
The last term is the pure transverse term.
The $\cos 3\theta$ is the lowest angular harmonic that is compatible with the $S_3$ symmetry of the teacher vectors: we must have $H(t, \theta + 2\pi/3) = H(t, \theta)$.

So, we can write the change in free energy as
\begin{align*}
    \Delta F
    = b \left[\lambda - \beta H(u) \right] + \frac{b^2}{2}
\end{align*}
For $0 < b \ll 1$, adding the committee will be favored when maximizing $H(u)$, as in the two-feature case.
Examining the formula \ref{eq:three-feature-H}, we see that $H$ is maximized for $\cos 3\theta = 1$, which gives three equivalent directions, corresponding to $e_1$, $e_2$, and $e_3$.

Then we can check when the various solutions are favored.
$t = 1$ corresponds to the direction of the committee being completely shared, giving $H_\text{sh} = 3c^3$.
We can check when this solution is locally stable.
Let $q = \sqrt{1-t^2} \ll 1$.
Then,
\begin{align*}
    H
    &= 3c^3 (1-q^2)^{3/2}
    + \frac{9}{2} cs^2 (1-q^2)q^2
    + \frac{3}{4} s^3 q^3 \cos (3\theta) \\
    &= 3c^3 + \frac{9}{2} c(s^2 - c^2)q^2
    + \frac{3}{4} s^3 q^3 \cos 3\theta + O(q^4)
\end{align*}
The $q^2$ term is negative (for recall we're trying to \emph{maximize} $H$), and this solution is stable, when $s^2 < c^2$, which is exactly when $\omega > 1/4$.

However, it turns out this is not necessarily the globally-stable point for $\omega > 1/4$.
We also need to consider the asymmetric solutions.
WLOG, let
\begin{align*}
(e_1^\top u, e_2^\top u, e_3^\top u)
= (x, y, y), \qquad x > y
\end{align*}
Then we want to maximize $H(u) = x^3 + 2y^3$.
To globally solve this constrained maximization problem, it helps to change coordinate systems.
Let
\begin{align}
    u = \sum_b \alpha_b e_b
\end{align}
If $G_{ab} = e_a^T e_b$ is the Gram matrix, we can write the projection of $u$ onto each feature vector as
\begin{align*}
    r_a  = e_a^\top u
    = \sum_b G_{ab} \alpha_b,
    \qquad \text{or }
    r &= G \alpha \\
    \implies
    \alpha &= G^{-1} r
\end{align*}
Thus, we can write the constraint that $u$ is a unit vector as
\begin{align*}
    |u|^2 = \alpha^\top G \alpha = r^\top G^{-1}  r = 1
\end{align*}
In this case with equally-correlated features,
\begin{align*}
    G
    = \begin{pmatrix}
        1 & \omega & \omega \\
        \omega & 1 & \omega \\
        \omega & \omega & 1
    \end{pmatrix}
    &= (1-\omega)I + \omega 1 1^\top \\
    G^{-1}
    &= \frac{1}{1-\omega} \left[I - \frac{\omega}{1 + 2\omega} 1 1^\top \right]
\end{align*}
We can now write the constrained maximization problem using a Lagrange mulitiplier $\mu$ as
\begin{align*}
    \mathcal{L} = \sum_a r_a^3 - \mu (r^\top G^{-1} r - 1)
\end{align*}
At the max, we have $\partial \mathcal{L} / \partial r_a = 0$, which gives
\begin{align}\label{eq:H-stationarity-lagrange}
    3r_a^2 = 2\mu (G^{-1} r)_a
\end{align}
We write the expression on the RHS as
\begin{align*}
    (G^{-1} r)_a
    = \frac{1}{1-\omega} \left[ r_a - \frac{\omega}{1 + 2\omega} (x + 2y) \right]
\end{align*}
Explicitly, the first and the symmetric components are
\begin{align*}
    (G^{-1} r)_1
    &= \frac{1}{1-\omega} \left[ x - \frac{\omega}{1 + 2\omega} (x + 2y) \right] \\
    &= \frac{(1 + \omega)x - 2\omega y}{(1-\omega)(1 + 2\omega)} \\
    (G^{-1}r)_{2,3}
    &= \frac{1}{1-\omega} \left[ y - \frac{\omega}{1 + 2\omega} (x + 2y) \right] \\
    &= \frac{y - \omega x}{(1-\omega)(1+2\omega)}
\end{align*}
Plugging these into equation \ref{eq:H-stationarity-lagrange}, we get the stationarity equations
\begin{align*}
    3x^2
    &= 2\mu \frac{(1 + \omega)x - 2\omega y}{(1-\omega)(1 + 2\omega)} \\
    3y^2
    &= 2 \mu \frac{y - \omega x}{(1-\omega)(1+2\omega)}
\end{align*}
Dividing and letting $R = x/y$, we have
\begin{align*}
    R^2
    &= \frac{(1+\omega)R - 2\omega}{1 - \omega R} \\
    \implies
    0 &= (R - 1) \left[\omega R^2 - (1 - \omega)R + 2\omega \right]
\end{align*}
$R=1$ is the $x = y$ symmetric case, but we're interested in the asymmetric one, so we can divide that out, giving the final stationarity equation
\begin{align}\label{eq:stationarity}
    0 &= \omega R^2 - (1-\omega)R + 2\omega \\
    \implies
    \omega &= \frac{R}{R^2 + R + 2}
\end{align}
This equation is necessary for an asymmetric stationary point; we still need to compare its value of $H$ with the symmetric solution.

Recall that for the symmetric case, $H_\text{sym}^2 = 3^2c^6 = 1/3 \times (1+2\omega)^3$, so
\begin{align*}
    \left( \frac{H_\text{asym}}{H_\text{sym}} \right)^2
    &= \frac{3y^6(R^3+2)^2}{(1+2\omega)^3}
\end{align*}
Then we need to replace $y$ and $\omega$ with $R$.
We can write $r = y(R, 1, 1)$, which lets us rewrite the $|u|^2 = 1$ constraint as
\begin{align*}
    1 &= \frac{y^2}{1-\omega} \left[ R^2 + 2 - \frac{\omega}{1+2\omega} (R+2)^2 \right] \\
    \implies y^2
    &= \frac{(1-\omega)(1+2\omega)}{(1+\omega)R^2 - 4\omega R +2}
\end{align*}
We'll let the denominator $D = (1+\omega)R^2 - 4\omega R + 2$, so
\begin{align*}
    y^6
    &= \frac{(1-\omega)^3 (1+2\omega)^3}{D^3} \\
    \implies
    \left( \frac{H_\text{asym}}{H_\text{sh}} \right)^2
    &= \frac{3(1-\omega)^3 (R^3+2)^2}{D^3}
\end{align*}
Then using equation \ref{eq:stationarity},
\begin{align*}
    1-\omega
    &= \frac{R^2+2}{R^2+R+2} \\
    D
    &= \frac{(R+2)(R^3+2)}{R^2+R+2}
\end{align*}
Plugging these into the above, we get
\begin{align*}
    \left( \frac{H_\text{asym}}{H_\text{sh}} \right)^2
    &= \frac{3(R^2+2)^3}{(R+2)^3(R^3+2)}
\end{align*}
The crossover point in favorability here is when the numerator and denominator are equal, which we can find:
\begin{align*}
    0
    &= 3(R^2+2)^3 - (R+2)^3(R^3+2)
    = 2(R-1)^3 (R^3-4)
\end{align*}
So, they are equally favored when (trivially) $R=1$, and nontrivially when $R = 2^{2/3}$.
Plugging this into equation \ref{eq:stationarity} gives
\begin{align*}
    \omega_c
    &= \frac{2^{2/3}}{2^{4/3} + 2^{2/3} + 2}
    = 2^{1/3} - 1
    \approx 0.259921 > \frac{1}{4}
\end{align*}
How large can $\omega$ get and keep the asymmetric solution locally stable?
Maximize equation \ref{eq:stationarity}:
\begin{align*}
    \frac{d\omega}{dR}
    &= \frac{(R^2 + R + 2) - R(2R+1)}{(R^2+R+2)^2} \\
    &= \frac{2-R^2}{(R^2+R+2)^2}
\end{align*}
So the max is at $R = 2^{1/2}$, and
\begin{align*}
    \omega_\text{spin}
    &= \frac{\sqrt{2}}{4+\sqrt{2}}
    \approx 0.261204
\end{align*}
So for correlations above $\omega_\text{spin}$, the shared direction is the only positive-valued local maximum of $H$.
As we lower $\omega$ through $\omega_\text{spin}$, an asymmetric local maximum appears, but the shared direction remains globally favored until $\omega_c = 2^{1/3}-1$.
At $\omega_c$, the two have equal $H$.
Between $\omega_c$ and $\omega = 1/4$, the symmetric solution remains locally stable, but the asymmetric solution is the global maximum.
Below $\omega = 1/4$, the symmetric solution loses local stability.
These statements concern which directions are favored at first recruitment; later population transitions require the overlap with the residual target.

\section{Spherical Harmonics for Any Number of Cubic Features}
\label{sec:cubic-spherical-harmonics}

The angular structure of $H$ in the three-feature calculation has a general explanation.
For $k$ unit teacher directions, the same Hermite overlap identity gives
\begin{align*}
    f_*(\mathbf x) &= B\sum_{\ell=1}^k \psi_3(\mathbf e_\ell^\top\mathbf x), \\
    \mathbb E_{\mathbf x}\left[f_*(\mathbf x)\psi_3(\mathbf u^\top\mathbf x)\right]
    &= B H(\mathbf u),
    \qquad H(\mathbf u) = \sum_{\ell=1}^k (\mathbf e_\ell^\top\mathbf u)^3.
\end{align*}
Here, $\mathbf u$ is a unit direction in the teacher span $\mathcal V$.
Thus, adding a committee to the unaligned state still changes the cost by
\begin{align*}
    F(b,\mathbf u)-F(0) = b[\lambda-BH(\mathbf u)]+\frac{b^2}{2}.
\end{align*}
We now allow arbitrary correlations between the teacher directions, rather than requiring all pairwise overlaps to equal $\omega$.
Let $m=\dim\mathcal V$ and use orthonormal coordinates in this span.
The angular variable then lies on $S^{m-1}$; for linearly independent teachers, $m=k$.
This distinction matters because the raw teacher overlaps describe an ellipsoid when the teachers are correlated, not a Euclidean unit sphere.

A generalized spherical harmonic of degree $j$ is the restriction to the unit sphere of a homogeneous degree-$j$ polynomial whose Euclidean Laplacian vanishes.
We will show that $H$ contains only degrees $1$ and $3$.
Write these two parts as $H^{(1)}$ and $H^{(3)}$, to distinguish them from the Hermite polynomial $H_3$ in the definition of $\psi_3$.

To find the decomposition, first extend $H$ off the sphere:
\begin{align*}
    P(\mathbf z) = \sum_{\ell=1}^k (\mathbf e_\ell^\top\mathbf z)^3,
    \qquad \mathbf v = \sum_{\ell=1}^k \mathbf e_\ell,
    \qquad \mathbf z\in\mathcal V.
\end{align*}
This polynomial is homogeneous of degree three, but it is generally not harmonic.
Since every teacher direction has unit norm,
\begin{align*}
    \Delta_{\mathbf z}(\mathbf e_\ell^\top\mathbf z)^3
    &= 6\|\mathbf e_\ell\|^2(\mathbf e_\ell^\top\mathbf z)
    = 6\mathbf e_\ell^\top\mathbf z, \\
    \Delta_{\mathbf z}P(\mathbf z) &= 6\mathbf v^\top\mathbf z.
\end{align*}
Here $\Delta_{\mathbf z}$ denotes the Laplacian in the $m$ orthonormal coordinates of $\mathcal V$, not the sample-scaling parameter $\Delta$.
The term we need to subtract follows from the product rule:
\begin{align*}
    \Delta_{\mathbf z}\left(\|\mathbf z\|^2\mathbf v^\top\mathbf z\right)
    &= 2m\mathbf v^\top\mathbf z+4\mathbf v^\top\mathbf z \\
    &= 2(m+2)\mathbf v^\top\mathbf z.
\end{align*}
Therefore, the polynomials
\begin{align*}
    Q_1(\mathbf z) &= \frac{3}{m+2}\mathbf v^\top\mathbf z, \\
    Q_3(\mathbf z) &= P(\mathbf z)-\|\mathbf z\|^2Q_1(\mathbf z)
\end{align*}
are homogeneous of degrees one and three, respectively, and both have zero Laplacian.
Restricting to $\|\mathbf u\|=1$ gives the desired decomposition:
\begin{align}
    H(\mathbf u) &= H^{(1)}(\mathbf u)+H^{(3)}(\mathbf u),
    \label{eq:H-harmonic-decomposition} \\
    H^{(1)}(\mathbf u) &= \frac{3}{m+2}\mathbf v^\top\mathbf u,
    \label{eq:H-degree-one} \\
    H^{(3)}(\mathbf u)
    &= \sum_{\ell=1}^k\left[(\mathbf e_\ell^\top\mathbf u)^3
    -\frac{3}{m+2}\mathbf e_\ell^\top\mathbf u\right].
    \label{eq:H-degree-three}
\end{align}
For $m\geq2$, these are spherical harmonics of degrees one and three on $S^{m-1}$.
Equivalently, expanding $H$ in any complete spherical-harmonic basis produces coefficients only at these two degrees.
The decomposition is exact, regardless of the correlations or the number of teachers.
If the span is one-dimensional, the sphere consists of just two points and $P(\mathbf z)=\|\mathbf z\|^2Q_1(\mathbf z)$, so $Q_3$ vanishes.

We can also express the answer using harmonics centered on each teacher direction.
For $m\geq2$, define
\begin{align*}
    Z_1^{(m)}(t)=t,
    \qquad Z_3^{(m)}(t)=\frac{(m+2)t^3-3t}{m-1}.
\end{align*}
The functions $Z_j^{(m)}(\mathbf e_\ell^\top\mathbf u)$ are zonal spherical harmonics: they depend only on the angle from $\mathbf e_\ell$.
Their normalization is $Z_j^{(m)}(1)=1$.
For $m>2$, they are the normalized Gegenbauer polynomials with parameter $(m-2)/2$; for $m=2$, they satisfy $Z_j^{(2)}(\cos\theta)=\cos(j\theta)$ for $j=1,3$.
Rearranging their definitions gives
\begin{align}
    H(\mathbf u)
    = \frac{3}{m+2}\sum_{\ell=1}^k Z_1^{(m)}(\mathbf e_\ell^\top\mathbf u)
    +\frac{m-1}{m+2}\sum_{\ell=1}^k Z_3^{(m)}(\mathbf e_\ell^\top\mathbf u).
    \label{eq:H-zonal-harmonics}
\end{align}

For the three-feature geometry above, $m=3$ when $\omega<1$, and $\mathbf v=3ce_0$.
The degree-one part is therefore $H^{(1)}(t,\theta)=9ct/5$.
Using $c^2+s^2=1$, equation \eqref{eq:three-feature-H} can be rewritten as
\begin{align*}
    H(t,\theta)
    &= \frac{9c}{5}t
    +\frac{3c}{5}(5c^2-3)P_3(t)
    +\frac{3}{4}s^3(1-t^2)^{3/2}\cos(3\theta), \\
    P_3(t)&=\frac{5t^3-3t}{2}.
\end{align*}
Here $P_3$ is the third Legendre polynomial.
The first term is a degree-one spherical harmonic, while the last two are degree-three spherical harmonics.
Thus, the $\cos(3\theta)$ term found earlier is one part of the full degree-three angular dependence.

The same proof also allows unequal feature strengths: replacing each summand in $H$ by $\eta_\ell(\mathbf e_\ell^\top\mathbf u)^3$ replaces $\mathbf v$ by $\sum_\ell\eta_\ell\mathbf e_\ell$ and weights each summand in \eqref{eq:H-degree-three} by $\eta_\ell$.
The degree-one component vanishes exactly when this weighted sum of directions is zero.
Although the target consists of cubic Hermite features in the Gaussian input, its overlap with a candidate committee generally contains both degree-one and degree-three spherical harmonics as a function of the committee direction.

\section{\texorpdfstring{$k$}{k} Features}

It turns out that generalizing the solution to $k$ features is not terribly difficult!
It will turn out that, for $k$ features, having some kind of symmetries in the problem will be important for deriving exact solutions.
However, when we reformulate into a certain kind of generalization of an eigenvalue problem called a $Z$-eigenvalue problem, it will become clearer how to attack even the fully asymmetric cases numerically.

We'll first need to generalize the stationarity condition into a tensorial form.
Start with the expression for $H(u)$ for $k$ features:
\begin{align*}
    H(u)
    &= \sum_{i=1}^k (e_i^\top u)^3,
    \qquad |u|^2 = 1
\end{align*}
We use the same notation as in the previous section, where $r_i = e_i^\top u$ is the overlap between $u$ and $e_i$, defining $r \in \mathbb{R}^k$ the overlap vector, and the Gram matrix $G \in \mathbb{R}^{k\times k}$ with $G_{ij} = e_i^\top e_j$.
Also, we will need a notation for elementwise squaring:
\begin{align*}
    r^{\circ 2} = (r_1^2, \dots,r_k^2)^\top
\end{align*}
We need to find the stationary points of $H$ subject to $|u| = 1$.
As before, introduce a Lagrange multiplier $\mu$:
\begin{align*}
    \mathcal{L}(u, \mu)
    &= H(u) - \frac{\mu}{2} (u^\top u - 1) \\
    \nabla_u H(u)
    &= 3 \sum_{i=1}^k (e_i^\top u)^2 e_i
    = 3 \sum_{i=1}^k r_i^2 e_i
\end{align*}
So stationarity requires
\begin{align*}
    3 \sum_{i=1}^k r_i^2 e_i
    = \mu u
\end{align*}
Dot both sides with $u$, and then using $|u|^2 = 1$, we get
\begin{align*}
    \mu
    &= 
    3 \sum_{i=1}^k r_i^3
    = 3 H(u) \\
    \implies
    \sum_i r_i^2 e_i
    &= H(u) u
\end{align*}
Then dot both sides with each teacher direction to get
\begin{align*}
    \sum_{i=1}^k (e_j^\top e_i) r_i^2
    &= H(u) (e_j^\top u) \\
    \implies
    \sum_{i}
    G_{ji} r_i^2 = H(u) r_j
\end{align*}
Or, written as vectors,
\begin{align*}
    \boxed{G r^{\circ 2} = H(u) r}
\end{align*}
This is the stationarity condition for any direction $u$; local stability additionally requires checking the curvature along the sphere.
We can see that this is very similar to an eigenvalue equation, except that the $G$ acts on $r$ \emph{quadratically}, rather than linearly.

We can write this as a tensorial equation.
Write $H(u)$ in the following way:
\begin{align*}
    H(u)
    &= \sum_{i=1}^k (e_i^\top u)^3 \\
    &= \sum_{i=1}^k \sum_{a,b,c=1}^d e_{ia} e_{ib} e_{ic} u_a u_b u_c \\
    &= \sum_{a,b,c=1}^d
    \underbrace{\left(\sum_{i=1}^k e_{ia} e_{ib} e_{ic}\right)}_{T_{abc}} 
    u_a u_b u_c
\end{align*}
Using Einstein summation notation, we can then write
\begin{align*}
    H(u) = T_{abc} u_a u_b u_c
\end{align*}
Where $T$ is a $d \times d \times d$ fully-symmetric tensor.

If you plug this into the stationarity equation above, you can find the following equivalent expression:
\begin{align*}
    \boxed{T_{abc} u_b u_c = H(u) u_a}
\end{align*}
This is the standard form for a $Z$-eigenvalue problem, with $\lambda_T = H(u)$.

This is also another way of seeing where the representation of $H$ as a sum of spherical harmonics comes from.
A rank three symmetric tensor breaks down in representation theory to a trace part, which has one index and is therefore a vector, and a traceless part.
In $k=3$, the traceless part is the degree-3 spherical harmonic, and the trace is the degree-1 harmonic.

Now say we're in the symmetric case, with $k\geq2$ and $0<\omega<1$.
Then
\begin{align*}
    G_{ij} = (1-\omega)\delta_{ij} + \omega
\end{align*}
So the components of the stationarity equation can be written
\begin{align*}
    (1-\omega)r_j^2
    + \omega \sum_{i=1}^k r_i^2
    &= H(u) r_j \\
    \implies
    (1-\omega)r_j^2
    - H(u) r_j
    + \omega \sum_{i=1}^k r_i^2 
    &= 0
\end{align*}
Take the above equation for two components, $r_j$ and $r_\ell$.
Subtracting gives
\begin{align*}
    (1-\omega)(r_j^2 - r_\ell^2)
    - H(u) (r_j - r_\ell)
    &= 0 \\
    \implies
    (r_j - r_\ell)
    [(1-\omega)(r_j + r_\ell) - H(u)]
    &= 0
\end{align*}
Thus, for any two components, either $r_j = r_\ell$, or
\begin{align*}
    r_j + r_\ell
    &= \frac{H(u)}{1-\omega}
\end{align*}
Now, suppose we have three pairwise distinct components $r_a$, $r_b$, and $r_c$.
Then
\begin{align*}
    (r_a + r_b)
    - (r_b + r_c)
    &= r_a - r_c = 0
\end{align*}
Thus, \emph{the components of $r$ can take only two distinct values}.

Label these $x$ and $y$.
By Vieta's root formulas (I hadn't heard of them either), we have
\begin{align}\label{eq:vieta}
    x + y = \frac{H}{1-\omega},
    \qquad
    xy = \frac{\omega r^2}{1-\omega}
\end{align}
where we defined $r^2 = \sum_i r_i^2$.
For positive correlation and positive $H$, $x+y > 0$, so at least one of $x$ and $y$ is positive, and $xy > 0$, so both have the same sign.
Thus, $x,y > 0$.

Let $x > y$.
We will next show that, if the configuration is a local maximum, then exactly one of the $r_i$'s is equal to $x$, and the rest are $y$.

Suppose there are at least two copies of the larger overlap $x$.
Then at least two teachers, WLOG $e_1$ and $e_2$, have overlap $x$: $e_1^\top u = e_2^\top u = x$.
Consider the normalized direction
\begin{align*}
    v
    = \frac{e_1 - e_2}{\sqrt{2(1-\omega)}}
\end{align*}
Moving along $v$ increases overlap with $e_1$ and decreases overlap with $e_2$, while not affecting other overlaps to first order.

$v$ is orthogonal to $u$:
\begin{align*}
    v^\top u
    = \frac{e_1^\top u - e_2^\top u}{\sqrt{2(1-\omega)}}
    = 0
\end{align*}
Thus, we can perturb $u$ while maintaining the normalization constraint by taking
\begin{align*}
    u(\varepsilon)
    &= u \cos\varepsilon + v \sin\varepsilon \\
    |u(\varepsilon)|^2
    &= |u|^2 \cos^2 \varepsilon
    + |v|^2 \sin^2 \varepsilon
    = 1
\end{align*}
The overlaps of $v$ with all teacher vectors are
\begin{align*}
    e_1^\top v
    &= \frac{1 - \omega}{\sqrt{2(1-\omega)}}
    = \sqrt{\frac{1-\omega}{2}}
    \equiv a \\
    e_2^\top v
    &= -a \\
    e_j^\top v &= 0,
    \qquad j \geq 3
\end{align*}
Thus, along the perturbation, the overlaps are
\begin{align*}
    r_1(\varepsilon)
    &= x \cos\varepsilon + a \sin\varepsilon \\
    r_2(\varepsilon)
    &= x \cos\varepsilon - a \sin\varepsilon \\
    r_j(\varepsilon)
    &= r_j \cos\varepsilon,
    \qquad j \geq 3
\end{align*}
So we are splitting the two equal $x$ overlaps in opposite directions, while keeping the rest unchanged apart from the $\cos\varepsilon$ normalization.

Now we'll calculate $H(u(\varepsilon))$.
\begin{align*}
    H(u(\varepsilon))
    &= (x\cos\varepsilon + a\sin\varepsilon)^3
    + (x\cos\varepsilon - a\sin\varepsilon)^3
    + \sum_{j=3}^k (r_j \cos\varepsilon)^3
\end{align*}
Using the identity
\begin{align*}
    (A+B)^3 + (A-B)^3
    = 2A^3 + 6AB^2
\end{align*}
we get
\begin{align*}
    H(u(\varepsilon))
    &= H(u) \cos^3\varepsilon
    + 6xa^2 \cos\varepsilon \sin^2\varepsilon \\
    &= H(u) \cos^3\varepsilon
    + 3x(1-\omega) \cos\varepsilon \sin^2\varepsilon \\
    &= H(u)
    + \left[3x(1-\omega) - \frac{3}{2} H(u) \right]\varepsilon^2 + O(\varepsilon^4)
\end{align*}
Using Vieta's root relations, equation \ref{eq:vieta}, we substitute $H(u) = (1-\omega)(x+y)$:
\begin{align*}
    3x(1-\omega) - \frac{3}{2} H(u)
    &= 3x(1-\omega)
    - \frac{3}{2}(1-\omega)(x+y) \\
    &= \frac{3}{2} (1-\omega)(x-y) \\
    \implies
    H(u(\varepsilon))
    &= H(u) + \frac{3}{2}(1-\omega)(x-y)\varepsilon^2 + O(\varepsilon^4)
\end{align*}
Since $\omega \in (0, 1)$ and $x > y$,
\begin{align*}
    \left. \frac{d^2 H(u(\varepsilon))}{d\varepsilon^2} \right|_{\varepsilon=0}
    = 3(1-\omega)(x-y) > 0
\end{align*}
Thus, we have constructed a direction in which $H$ curves upward, meaning the original stationary point could not have been a local maximum.

So, we have proved:
\begin{enumerate}
    \item Every positive-valued stationary point for $\omega \in (0, 1)$ has at most two overlap values $x \geq y > 0$.
    \item If $x > y$, any configuration containing two or more copies of $x$ has a tangent direction with positive curvature.
    \item Therefore, every local maximum with two distinct overlap values must have exactly one copy of $x$.
\end{enumerate}
Since $H(e_i)=1+(k-1)\omega^3>0$, every global maximizer has positive $H$ and is covered by this argument.
Hence, every global maximizer is one of two forms:
\begin{align*}
    r = (x, x, \dots, x)
    \quad \text{or} \quad
    r = (x, y, \dots, y),
    \qquad x > y
\end{align*}
The first case is the fully-shared, symmetric solution, and the second is the specialized, symmetry-broken solution.

Now, we want to derive the critical $\omega$ value for the asymmetric solution to be the globally stable one.
Take, WLOG, $e_1^\top u = x$, and all the other overlaps to be $y$.
Also let $R = x/y > 1$.
Then we want to maximize
\begin{align*}
    H(u)
    &= \sum_{i=1}^k (e_i^\top u)^3,
    \qquad |u|^2 = 1
\end{align*}
The Lagrange multiplier equation is
\begin{align*}
    3 \sum_i (e_i^\top u)^2 e_i
    = \mu u
    \implies
    3 H(u) = \mu
    \implies \sum_i (e_i^\top u)^2 e_i = H(u)u
\end{align*}
Taking inner products with $e_1$ and any other teacher direction gives, respectively,
\begin{align*}
    H(u) x
    &= x^2 + \omega(k-1)y^2 \\
    H(u) y
    &= \omega x^2 + [1 + (k-2)\omega]y^2
\end{align*}
Substitute equation \ref{eq:vieta} in for $H(u)$ in the second equation here to get
\begin{align*}
    (1-\omega)(x+y)y
    &= \omega x^2 + [1 + (k-2)\omega]y^2 \\
    \omega x^2 - (1-\omega)xy + \omega(k-1)y^2 &= 0
\end{align*}
Then divide by $y^2$ and use $R = x/y$:
\begin{align}\label{eq:k-feature-stationarity}
    \omega R^2 - (1-\omega)R + \omega(k-1) &= 0 \\
    \implies
    \omega
    &= \frac{R}{R^2 + R + k-1}
\end{align}
We derived a similar equation in the $k=3$ case above.

Again, similarly to above, we want to get the ratio of the $H_\text{asym}$ and $H_\text{sym}$ values.
For the symmetric direction, normalize the sum of the teacher vectors:
\begin{align*}
    e_0
    &= \frac{\sum_{i=1}^k e_i}{\sqrt{k[1+(k-1)\omega]}}, \\
    e_i^\top e_0
    &= \sqrt{\frac{1+(k-1)\omega}{k}}
    \equiv c_k.
\end{align*}
All $k$ overlaps are equal to $c_k$, so
\begin{align*}
    H_\text{sym} &= kc_k^3,
    \qquad H_\text{sym}^2=\frac{[1+(k-1)\omega]^3}{k}.
\end{align*}
For the asymmetric direction, $r=y(R,1,\ldots,1)$ gives
\begin{align*}
    H_\text{asym}=y^3(R^3+k-1).
\end{align*}
Thus the ratio we need is
\begin{align*}
    \left(\frac{H_\text{asym}}{H_\text{sym}}\right)^2
    =\frac{k y^6(R^3+k-1)^2}{[1+(k-1)\omega]^3}.
\end{align*}
As in the three-feature case, the remaining step is to eliminate $y$ using the normalization constraint, and then eliminate $\omega$ using the stationarity equation.

The inverse Gram matrix for $k$ equally correlated teachers is
\begin{align*}
    G^{-1}
    =\frac{1}{1-\omega}
    \left[I-\frac{\omega}{1+(k-1)\omega}\mathbf 1\mathbf 1^\top\right],
\end{align*}
where $\mathbf 1$ is the vector of $k$ ones.
Since $|u|^2=r^\top G^{-1}r=1$, we get
\begin{align*}
    1
    &=\frac{y^2}{1-\omega}
    \left[R^2+k-1-\frac{\omega(R+k-1)^2}{1+(k-1)\omega}\right] \\
    \implies y^2
    &=\frac{(1-\omega)[1+(k-1)\omega]}{D},
\end{align*}
where, expanding the denominator,
\begin{align*}
    D
    &=[1+(k-1)\omega](R^2+k-1)-\omega(R+k-1)^2 \\
    &=[1+(k-2)\omega]R^2-2(k-1)\omega R+k-1.
\end{align*}
Cubing $y^2$ and substituting into the ratio cancels the factor $[1+(k-1)\omega]^3$:
\begin{align*}
    \left(\frac{H_\text{asym}}{H_\text{sym}}\right)^2
    =\frac{k(1-\omega)^3(R^3+k-1)^2}{D^3}.
\end{align*}

Now use equation \eqref{eq:k-feature-stationarity} to replace $\omega$ by $R$.
The two factors we need simplify to
\begin{align*}
    1-\omega
    &=\frac{R^2+k-1}{R^2+R+k-1}, \\
    D
    &=\frac{R^4+(k-1)R^3+(k-1)R+(k-1)^2}{R^2+R+k-1} \\
    &=\frac{(R+k-1)(R^3+k-1)}{R^2+R+k-1}.
\end{align*}
Substituting these expressions and cancelling common factors gives
\begin{align}\label{eq:k-feature-H-ratio}
    \boxed{\left(\frac{H_\text{asym}}{H_\text{sym}}\right)^2
    =\frac{k(R^2+k-1)^3}{(R+k-1)^3(R^3+k-1)}.}
\end{align}
Both values of $H$ are positive, so comparing their squares gives the same ordering as comparing $H$ itself.

The crossover occurs when the numerator and denominator of \eqref{eq:k-feature-H-ratio} are equal.
To see the factorization more clearly, write $q=k-1$ temporarily:
\begin{align*}
    0
    &=(q+1)(R^2+q)^3-(R+q)^3(R^3+q) \\
    &=q\left[R^6-3R^5+3R^4-(1+q^2)R^3
       +3q^2R^2-3q^2R+q^2\right] \\
    &=q(R-1)^3(R^3-q^2).
\end{align*}
The root $R=1$ is the symmetric direction itself.
For $k>2$, the nontrivial crossing is therefore
\begin{align*}
    R_c^3=(k-1)^2,
    \qquad R_c=(k-1)^{2/3}>1.
\end{align*}
Moreover, since $R>1$ and all denominators in \eqref{eq:k-feature-H-ratio} are positive,
\begin{align*}
    H_\text{asym}>H_\text{sym}
    \quad\Longleftrightarrow\quad R>R_c.
\end{align*}

We still need to check which stationary branch this crossing lies on.
Solving the quadratic in equation \eqref{eq:k-feature-stationarity} gives
\begin{align*}
    R_\pm(\omega)
    =\frac{1-\omega\pm\sqrt{(1-\omega)^2-4(k-1)\omega^2}}{2\omega}.
\end{align*}
The roots are real only for
\begin{align*}
    \omega\leq\frac{1}{1+2\sqrt{k-1}}
    \equiv\omega_\text{spin}(k),
\end{align*}
and their product is $R_+R_-=k-1$.
Thus $R_-\leq\sqrt{k-1}\leq R_+$.
For $k>2$, $R_c=(k-1)^{2/3}>\sqrt{k-1}$, so the smaller root, whenever it satisfies $R_->1$, always has lower $H$ than the shared direction.
Only $R_+$ can be globally favored.
On this larger-root branch,
\begin{align*}
    \frac{d\omega}{dR}
    &=\frac{(R^2+R+k-1)-R(2R+1)}{(R^2+R+k-1)^2} \\
    &=\frac{k-1-R^2}{(R^2+R+k-1)^2}<0,
    \qquad R>\sqrt{k-1}.
\end{align*}
So increasing $R$ past $R_c$ means decreasing $\omega$ below the critical correlation.
Since we have already exhausted the possible forms of a global maximizer, this comparison determines the global transition, rather than only a local-stability boundary.

Finally, substituting $R_c$ into equation \eqref{eq:k-feature-stationarity} gives
\begin{align}\label{eq:k-feature-omega-critical}
    \omega_c(k)
    &=\frac{(k-1)^{2/3}}{(k-1)^{4/3}+(k-1)^{2/3}+k-1} \\
    &=\boxed{\frac{1}{1+(k-1)^{1/3}+(k-1)^{2/3}}.}
\end{align}
For $0<\omega<\omega_c(k)$, one of the $k$ equivalent asymmetric directions is favored; for $\omega>\omega_c(k)$, the shared direction is favored.
At $\omega_c(k)$ and $k>2$, the shared direction and the $k$ distinct asymmetric directions have equal $H$.
The asymmetric overlap ratio is still $R_c>1$ there, so the maximizing direction changes discontinuously.

As a check, for three features this reduces to
\begin{align*}
    \omega_c(3)
    =\frac{1}{1+2^{1/3}+2^{2/3}}
    =2^{1/3}-1,
\end{align*}
exactly as above.
For $k=2$, $R_c=1$: the asymmetric branch merges into the shared direction continuously at $\omega_c(2)=1/3$, rather than crossing it at a distinct direction.
At $\omega=0$, the maximizers are the individual orthogonal teachers, obtained as the limit $R_+\to\infty$; at $\omega=1$, all teachers coincide and only the shared direction remains.
For large $k$, the critical correlation scales as $\omega_c(k)\sim k^{-2/3}$.
As in the three-feature calculation, this is the correlation transition in the direction first recruited from the unaligned state, not the later splitting threshold of an occupied committee.

\section{Local Stability with Constant Row Sums}
\label{sec:constant-row-sum-stability}

The shared direction can remain stationary even when the pairwise correlations are not all equal.
A useful condition is that every teacher has the same total overlap with the others, including itself:
\begin{align*}
    G\mathbf 1=s\mathbf 1,
    \qquad \sum_{j=1}^k G_{ij}=s>0.
\end{align*}
Here $\mathbf 1$ is the vector of $k$ ones, and $G_{ij}=e_i^\top e_j$ is still the Gram matrix of unit teacher directions.
We keep equal feature amplitudes, so the angular objective is
\begin{align*}
    H(u)=\sum_{i=1}^k(e_i^\top u)^3,
    \qquad |u|=1.
\end{align*}
We will first show that the normalized sum of the teachers is stationary, then determine when every small angular perturbation lowers $H$ to quadratic order.
We work in the teacher span $\mathcal V$ and assume $\dim\mathcal V\geq2$, so there are angular directions to perturb.

Start by defining $w=\sum_i e_i$.
The constant-row-sum condition fixes its norm:
\begin{align*}
    |w|^2
    &=\sum_{i,j=1}^k e_i^\top e_j
    =\sum_{i=1}^k\sum_{j=1}^k G_{ij}
    =ks.
\end{align*}
Thus the shared unit direction is
\begin{align*}
    u_0=\frac{\sum_i e_i}{\sqrt{ks}}.
\end{align*}
Its overlap with any teacher is
\begin{align*}
    e_i^\top u_0
    &=\frac{\sum_j e_i^\top e_j}{\sqrt{ks}}
    =\frac{s}{\sqrt{ks}}
    =\sqrt{\frac{s}{k}}
    \equiv c.
\end{align*}
All teachers therefore have the same overlap $c$, even though individual entries of $G$ can differ.
In particular,
\begin{align*}
    H(u_0)=kc^3=sc,
\end{align*}
where we used $kc^2=s$.

Now evaluate the gradient at this direction:
\begin{align*}
    \nabla_u H(u_0)
    &=3\sum_i(e_i^\top u_0)^2e_i
    =3c^2\sum_i e_i \\
    &=3c^2\sqrt{ks}\,u_0
    =3sc\,u_0
    =3H(u_0)u_0.
\end{align*}
This is precisely the stationarity equation from the previous section.
The gradient is radial, so it has no component along the sphere.
Stationarity alone does not tell us whether the shared direction is a maximum; we need its curvature along tangent directions.

Choose any unit vector $v\in\mathcal V$ perpendicular to $u_0$.
As in the earlier splitting calculation, use a great-circle perturbation:
\begin{align*}
    u(\varepsilon)=u_0\cos\varepsilon+v\sin\varepsilon,
    \qquad |v|=1,
    \qquad v^\top u_0=0.
\end{align*}
The constraint is preserved exactly because
\begin{align*}
    |u(\varepsilon)|^2
    &=|u_0|^2\cos^2\varepsilon+|v|^2\sin^2\varepsilon
      +2u_0^\top v\cos\varepsilon\sin\varepsilon \\
    &=\cos^2\varepsilon+\sin^2\varepsilon=1.
\end{align*}
Let $a_i=e_i^\top v$ be the overlaps of the perturbation direction with the teachers.
Their sum vanishes:
\begin{align*}
    \sum_i a_i
    =\left(\sum_i e_i\right)^\top v
    =\sqrt{ks}\,u_0^\top v=0.
\end{align*}
The perturbed teacher overlaps are $e_i^\top u(\varepsilon)=c\cos\varepsilon+a_i\sin\varepsilon$, so expanding each cube gives
\begin{align*}
    H(u(\varepsilon))
    &=kc^3\cos^3\varepsilon
      +3c^2\cos^2\varepsilon\sin\varepsilon\sum_i a_i \\
    &\quad+3c\cos\varepsilon\sin^2\varepsilon\sum_i a_i^2
      +\sin^3\varepsilon\sum_i a_i^3 \\
    &=H(u_0)\cos^3\varepsilon
      +3c\cos\varepsilon\sin^2\varepsilon\sum_i a_i^2
      +\sin^3\varepsilon\sum_i a_i^3.
\end{align*}
The term linear in the $a_i$ cancels because $\sum_i a_i=0$.
Using
\begin{align*}
    \cos^3\varepsilon&=1-\frac32\varepsilon^2+O(\varepsilon^4), \\
    \cos\varepsilon\sin^2\varepsilon&=\varepsilon^2+O(\varepsilon^4),
    \qquad \sin^3\varepsilon=\varepsilon^3+O(\varepsilon^5),
\end{align*}
we find
\begin{align*}
    H(u(\varepsilon))
    &=H(u_0)
      +\left[3c\sum_i a_i^2-\frac32 H(u_0)\right]\varepsilon^2
      +\left(\sum_i a_i^3\right)\varepsilon^3+O(\varepsilon^4) \\
    &=H(u_0)
      +\frac{3c}{2}\left[2\sum_i a_i^2-s\right]\varepsilon^2
      +\left(\sum_i a_i^3\right)\varepsilon^3+O(\varepsilon^4).
\end{align*}
Thus the curvature in direction $v$ is
\begin{align}\label{eq:row-sum-tangent-curvature}
    \left.\frac{d^2H(u(\varepsilon))}{d\varepsilon^2}\right|_{\varepsilon=0}
    =3c\left[2\sum_i(e_i^\top v)^2-s\right].
\end{align}
The negative term comes from maintaining the unit-norm constraint: tilting away from $u_0$ reduces the shared component by the factor $\cos\varepsilon$.

We can now turn the maximization over perturbations into a matrix eigenvalue problem.
Let $E$ be the matrix whose columns are the teacher vectors, so
\begin{align*}
    E=(e_1\ \cdots\ e_k),
    \qquad G=E^\top E,
    \qquad \sum_i(e_i^\top v)^2=|E^\top v|^2=v^\top EE^\top v.
\end{align*}
Equivalently, the Hessian of the constrained Lagrangian at $u_0$, restricted to tangent vectors, is
\begin{align*}
    \nabla_u^2\mathcal L(u_0,\mu)
    &=6c\,EE^\top-\mu I
    =3c(2EE^\top-sI),
    \qquad \mu=3H(u_0)=3sc.
\end{align*}
This agrees with the great-circle calculation in \eqref{eq:row-sum-tangent-curvature}.

To express the relevant eigenvalues using $G$, first note that $\mathbf 1^\perp$ is an invariant subspace of $G$.
Indeed, if $\mathbf 1^\top a=0$, then symmetry of $G$ gives
\begin{align*}
    \mathbf 1^\top Ga=(G\mathbf 1)^\top a=s\mathbf 1^\top a=0.
\end{align*}
Suppose $a$ is a unit eigenvector in this subspace, with positive eigenvalue $\gamma$:
\begin{align*}
    Ga=\gamma a,
    \qquad |a|=1,
    \qquad \mathbf 1^\top a=0.
\end{align*}
Define $v=Ea/\sqrt\gamma$.
Its norm and its overlap with the shared direction are
\begin{align*}
    |v|^2
    &=\frac{a^\top E^\top Ea}{\gamma}
    =\frac{a^\top Ga}{\gamma}=1, \\
    u_0^\top v
    &=\frac{(E^\top u_0)^\top a}{\sqrt\gamma}
    =\frac{c\mathbf 1^\top a}{\sqrt\gamma}=0.
\end{align*}
So $v$ is an allowed unit tangent direction.
Moreover,
\begin{align*}
    EE^\top v
    &=\frac{EE^\top Ea}{\sqrt\gamma}
    =\frac{EGa}{\sqrt\gamma}
    =\gamma v.
\end{align*}
Conversely, a positive-eigenvalue tangent eigenvector $v$ of $EE^\top$ maps back to $a=E^\top v/\sqrt\gamma$.
This $a$ has unit norm, satisfies $Ga=\gamma a$, and obeys $\mathbf 1^\top a=\sqrt{ks}\,u_0^\top v/\sqrt\gamma=0$.
Therefore the tangent eigenvalues of $EE^\top$ within the teacher span are exactly the positive eigenvalues of $G$ on $\mathbf 1^\perp$.
If the teachers are linearly dependent, zero eigenvalues of $G$ correspond to $Ea=0$ and add no physical perturbation directions.

The largest quadratic form on unit tangent vectors is consequently
\begin{align*}
    \max_{\substack{v\in\mathcal V,\ |v|=1\\v^\top u_0=0}}
    \sum_i(e_i^\top v)^2
    =\lambda_{\max}\!\left(G\big|_{\mathbf 1^\perp}\right).
\end{align*}
Here the restriction means that we exclude the all-ones direction and keep all eigenmodes orthogonal to it.
Since $c>0$, equation \eqref{eq:row-sum-tangent-curvature} is negative for every nonzero tangent direction precisely when
\begin{align}\label{eq:constant-row-sum-stability}
    \boxed{2\lambda_{\max}\!\left(G\big|_{\mathbf 1^\perp}\right)<s.}
\end{align}
This is the condition for a negative-definite tangent Hessian, and hence a strict local maximum of $H$.
If the inequality is reversed, an eigenvector with $2\gamma>s$ supplies a direction in which $H$ increases, so $u_0$ is not a local maximum.
If one also allows perturbations outside the teacher span, those have $E^\top v=0$ and curvature $-3cs<0$; they do not create an additional instability.

Equality requires more care.
When $2\lambda_{\max}=s$, the quadratic term vanishes along at least one mode, so the cubic and higher-order terms decide stability.
For example, if $\sum_i a_i^3\ne0$ along such a mode, one sign of $\varepsilon$ increases $H$ at cubic order.
If that cubic term also vanishes, we must continue the expansion.
Thus \eqref{eq:constant-row-sum-stability} is an exact test for negative quadratic curvature, but is not a necessary condition for every possible strict local maximum with a degenerate Hessian.

As a check, return to equal pairwise correlations:
\begin{align*}
    G=(1-\omega)I+\omega\mathbf 1\mathbf 1^\top,
    \qquad s=1+(k-1)\omega.
\end{align*}
For any $a\perp\mathbf 1$, the second term annihilates $a$, giving
\begin{align*}
    Ga=(1-\omega)a,
    \qquad \lambda_{\max}\!\left(G\big|_{\mathbf 1^\perp}\right)=1-\omega.
\end{align*}
The stability condition reduces to
\begin{align*}
    2(1-\omega)&<1+(k-1)\omega \\
    \implies 1&<(k+1)\omega \\
    \implies \omega&>\frac{1}{k+1}.
\end{align*}
For $k=3$, this is the $1/4$ local-stability boundary found earlier.
It is smaller than the global transition $\omega_c(3)=2^{1/3}-1$: the shared direction can resist every sufficiently small tilt while a more distant asymmetric direction has higher $H$.
For $k=2$ at the marginal value $\omega=1/3$, the great-circle expression instead gives
\begin{align*}
    \frac{H(u(\varepsilon))}{H(u_0)}
    =\frac32\cos\varepsilon-\frac12\cos^3\varepsilon
    =1-\frac38\varepsilon^4+O(\varepsilon^6).
\end{align*}
The shared direction is still a strict local maximum at this boundary, even though its quadratic curvature is zero.

\section{Exact Splitting of Orthogonal Feature Groups}
\label{sec:orthogonal-feature-groups}

There is one correlation structure for which we can separate the entire population problem, including after the first committees have appeared.
Suppose the teachers can be partitioned into groups $I_1,\ldots,I_A$, with
\begin{align*}
    \mathcal V_a&=\operatorname{span}\{e_i:i\in I_a\},
    \qquad \mathcal V=\bigoplus_{a=1}^A\mathcal V_a, \\
    e_i^\top e_j&=0
    \qquad (i\in I_a,\ j\in I_b,\ a\ne b).
\end{align*}
Thus the Gram matrix is block diagonal.
Correlations within each group can be arbitrary; we only require different groups to span mutually orthogonal subspaces.
We keep the Gaussian input law and the reduced objective
\begin{align*}
    F(\mathcal C)
    &=\lambda\sum_j b_j+\frac12\mathbb E[f_*(\mathbf x)-f_{\mathcal C}(\mathbf x)]^2, \\
    f_{\mathcal C}(\mathbf x)
    &=\sum_j b_j\psi_3(u_j^\top\mathbf x),
    \qquad b_j\geq0,\quad |u_j|=1,\quad u_j\in\mathcal V,
\end{align*}
with fixed $\lambda>0$ and a freely varying number of committees.
We will show that a committee spanning several orthogonal groups can always be replaced by separate committees at lower cost.

Let $P_a$ be the orthogonal projection onto $\mathcal V_a$, and write $\mathbf x_a=P_a\mathbf x$.
Since $\mathbf x\sim\mathcal N(0,I_N)$, the random vectors $\mathbf x_a$ are independent: their cross-covariances are $P_aP_b=0$ for $a\ne b$, and they are jointly Gaussian.
The teacher is already a sum of functions of these separate inputs:
\begin{align*}
    f_*(\mathbf x)&=\sum_{a=1}^A f_{*,a}(\mathbf x_a), \\
    f_{*,a}(\mathbf x_a)&=B\sum_{i\in I_a}\psi_3(e_i^\top\mathbf x_a).
\end{align*}
Each group contribution has zero mean.
The important question is whether a student committee mixing several groups could nevertheless improve the fit enough to justify its cost.

\subsection{Separating the first-recruited direction}

First consider the angular objective $H$.
Any unit direction in $\mathcal V$ can be decomposed as
\begin{align*}
    u=\sum_a t_av_a,\qquad
    t_a=|P_au|\geq0,\qquad
    v_a=\frac{P_au}{t_a}\in\mathcal V_a,\qquad
    \sum_a t_a^2=1.
\end{align*}
Terms with $t_a=0$ can be omitted.
The coefficients $t_a$ are nonnegative lengths; signs are included in $v_a$.
For a teacher in group $a$, orthogonality gives $e_i^\top u=t_ae_i^\top v_a$.
Therefore
\begin{align*}
    H(u)
    &=\sum_a\sum_{i\in I_a}(e_i^\top u)^3
      =\sum_a t_a^3H_a(v_a), \\
    H_a(v)&=\sum_{i\in I_a}(e_i^\top v)^3,
    \qquad M_a=\max_{\substack{v\in\mathcal V_a\\|v|=1}}H_a(v).
\end{align*}
Since $H_a(-v)=-H_a(v)$, each $M_a$ is nonnegative.
Also $0\leq t_a\leq1$, so $t_a^3\leq t_a^2$.
Putting $M=\max_a M_a$, we obtain
\begin{align*}
    H(u)
    &\leq\sum_a t_a^3M_a
    \leq M\sum_a t_a^3
    \leq M\sum_a t_a^2=M.
\end{align*}
The bound is attained by choosing $u$ entirely in a group whose maximum is $M$, along one of its maximizing directions.
Thus
\begin{align}\label{eq:orthogonal-H-maximum}
    \boxed{\max_{\substack{u\in\mathcal V\\|u|=1}}H(u)=\max_a M_a.}
\end{align}
For $M>0$, the first recruitment threshold is consequently $B_{\mathrm{on}}=\lambda/M$.
If several groups tie for the largest positive $M_a$, their maximizing directions tie at onset, but a direction mixing them does not: two or more nonzero $t_a$ give $\sum_a t_a^3<1$.

\subsection{What is removed when a committee is split}

The preceding argument only concerns onset.
To separate the full population optimization, we need to compare the actual student functions and their squared errors.
Take one committee with output $b\psi_3(u^\top\mathbf x)$ and use the same decomposition $u=\sum_a t_av_a$.
Conditioned on $\mathbf x_a$, put $z=v_a^\top\mathbf x_a$ and write
\begin{align*}
    u^\top\mathbf x=t_az+\sqrt{1-t_a^2}\,Z,
    \qquad Z\sim\mathcal N(0,1),
\end{align*}
where $Z$ is independent of $\mathbf x_a$.
For $t_a=1$ the second term is absent.
Using $\mathbb E Z=\mathbb E Z^3=0$ and $\mathbb E Z^2=1$, we find
\begin{align*}
    \mathbb E[\psi_3(u^\top\mathbf x)\mid\mathbf x_a]
    &=\frac{t_a^3z^3+3t_a(1-t_a^2)z-3t_az}{\sqrt6} \\
    &=\frac{t_a^3(z^3-3z)}{\sqrt6}
      =t_a^3\psi_3(v_a^\top\mathbf x_a).
\end{align*}
Thus the part depending on group $a$ alone is itself a cubic committee, with strength $bt_a^3$.
This suggests the replacement
\begin{align}\label{eq:orthogonal-committee-replacement}
    b\psi_3(u^\top\mathbf x)
    \quad\longrightarrow\quad
    b\sum_a t_a^3\psi_3(v_a^\top\mathbf x_a).
\end{align}
This replacement changes the function: it removes terms that involve inputs from more than one group.

For example, with two groups let $u=tv+sw$, where $t^2+s^2=1$ and $v,w$ are unit vectors in the respective subspaces.
Then $Z=v^\top\mathbf x$ and $W=w^\top\mathbf x$ are independent standard Gaussians.
Expanding the cubic and using $1-t^2=s^2$ and $1-s^2=t^2$ gives
\begin{align*}
    \psi_3(tZ+sW)
    &=t^3\psi_3(Z)+s^3\psi_3(W) \\
    &\quad+\frac{3t^2s}{\sqrt6}(Z^2-1)W
      +\frac{3ts^2}{\sqrt6}Z(W^2-1) \\
    &=t^3\psi_3(Z)+s^3\psi_3(W)
      +\sqrt3\,t^2s\,\psi_2(Z)W
      +\sqrt3\,ts^2\,Z\psi_2(W),
\end{align*}
where $\psi_2(z)=(z^2-1)/\sqrt2$.
The first two terms are retained by \eqref{eq:orthogonal-committee-replacement}.
The last two are mixed terms, and averaging either one over either input leaves zero.

For any number of groups, define the removed part by
\begin{align*}
    m_u(\mathbf x)
    =\psi_3(u^\top\mathbf x)
      -\sum_a t_a^3\psi_3(v_a^\top\mathbf x_a).
\end{align*}
The conditional expectation calculation gives
\begin{align*}
    \mathbb E[m_u\mid\mathbf x_a]
    &=t_a^3\psi_3(v_a^\top\mathbf x_a)
      -t_a^3\psi_3(v_a^\top\mathbf x_a) \\
    &\quad-\sum_{b\ne a}t_b^3\mathbb E[\psi_3(v_b^\top\mathbf x_b)]
      =0.
\end{align*}
Here independence lets us average each other group's term unconditionally, and every cubic Hermite has mean zero.
Consequently, for any square-integrable function $h_a$ of group $a$,
\begin{align*}
    \mathbb E[m_u(\mathbf x)h_a(\mathbf x_a)]
    =\mathbb E\!\left[h_a(\mathbf x_a)\mathbb E[m_u\mid\mathbf x_a]\right]
    =0.
\end{align*}
The removed part is therefore orthogonal, in the squared-error inner product, to every sum of functions of individual groups.

\subsection{Comparing the full free energies}

Now split every committee in a population $\mathcal C$.
Writing $u_j=\sum_a t_{ja}v_{ja}$, the original and split students are
\begin{align*}
    f_{\mathcal C}&=g+m, \\
    g(\mathbf x)&=\sum_{j,a}b_jt_{ja}^3\psi_3(v_{ja}^\top\mathbf x_a), \\
    m(\mathbf x)&=\sum_j b_jm_{u_j}(\mathbf x).
\end{align*}
Both $f_*$ and $g$ are sums of functions of individual groups.
The orthogonality just proved therefore implies $\mathbb E[(f_*-g)m]=0$.
Expanding the error gives the exact identity
\begin{align*}
    \mathbb E[f_*-f_{\mathcal C}]^2
    &=\mathbb E[(f_*-g)-m]^2 \\
    &=\mathbb E[f_*-g]^2-2\mathbb E[(f_*-g)m]+\mathbb E[m^2] \\
    &=\mathbb E[f_*-g]^2+\mathbb E[m^2].
\end{align*}
Thus splitting cannot increase the fitting error.
We use the total removed function $m$ here, so the argument also allows mixed terms from different committees to cancel one another.

The total strength of the split population satisfies
\begin{align*}
    \sum_{j,a}b_jt_{ja}^3
    \leq\sum_j b_j\sum_a t_{ja}^2
    =\sum_j b_j.
\end{align*}
Combining the two changes yields
\begin{align}\label{eq:orthogonal-free-energy-difference}
    \boxed{
    F(\mathcal C_{\mathrm{split}})-F(\mathcal C)
    =-\lambda\sum_j b_j\left(1-\sum_a t_{ja}^3\right)
      -\frac12\mathbb E[m^2]\leq0.}
\end{align}
If any committee of positive strength occupies more than one group, its sum $\sum_a t_{ja}^3$ is strictly less than one.
For $\lambda>0$, splitting then strictly lowers the free energy, even if the total mixed function happens to vanish.
It follows that every positive-strength committee in an optimal population lies entirely within one group.

Once committees have this form, write $f_{\mathcal C}=\sum_a f_{\mathcal C_a}(\mathbf x_a)$ and define the group residuals
\begin{align*}
    R_a(\mathbf x_a)=f_{*,a}(\mathbf x_a)-f_{\mathcal C_a}(\mathbf x_a).
\end{align*}
These residuals are centered and independent across groups, so for $a\ne b$,
\begin{align*}
    \mathbb E[R_aR_b]=\mathbb E[R_a]\mathbb E[R_b]=0.
\end{align*}
The objective therefore separates:
\begin{align*}
    F(\mathcal C)
    &=\sum_a\left[
       \lambda\sum_{j\in\mathcal C_a}b_j
       +\frac12\mathbb E R_a^2\right]
      \equiv\sum_a F_a(\mathcal C_a).
\end{align*}
Conversely, combining any collection of group populations gives an admissible population for the full problem.
Together with \eqref{eq:orthogonal-free-energy-difference}, this proves
\begin{align}\label{eq:orthogonal-population-minimum}
    \boxed{\inf_{\mathcal C}F(\mathcal C)
    =\sum_a\inf_{\mathcal C_a}F_a(\mathcal C_a).}
\end{align}
In particular, whenever group minimizers are chosen, their union is a global minimizer.
This holds at every signal strength for the reduced objective with fixed $\lambda$.
Each group can recruit and split its committees according to its own correlations, without changing the optimization in the other groups.
For a nonzero group overlap, its individual onset is $B_{\mathrm{on},a}=\lambda/M_a$; groups with smaller $M_a$ can start learning after the first group has already recruited.

As a simple check, if all teachers are mutually orthogonal, each group can contain just one teacher.
After absorbing signs into directions and combining committees along that teacher, the group objective is
\begin{align*}
    F_a(b_a)=\lambda b_a+\frac12(B-b_a)^2,\qquad b_a\geq0.
\end{align*}
Differentiating gives $F_a'(b_a)=\lambda+b_a-B$, so
\begin{align*}
    b_a=(B-\lambda)_+,\qquad (z)_+=\max\{z,0\}.
\end{align*}
Each teacher is learned by its own committee, with no advantage from a direction mixing teachers.

The assumptions matter for the exact separation.
The proof uses Gaussian independence and the cubic Hermite conditional expectation, together with the penalty on total strength and the freedom to introduce additional committees.
A fixed upper limit on their number, an extra cost per committee, or a constraint coupling group strengths would require a separate argument.
Likewise, orthogonal teacher directions alone do not give the required independence for a general non-Gaussian input law.

\appendix
\section{Generalized Spherical Harmonics}
\label{app:generalized-spherical-harmonics}

Ordinary spherical harmonics describe angular dependence on $S^2$, the unit sphere in $\R^3$.
Here, ``generalized'' means that we allow the ambient space to be $\R^m$ and the sphere to be
\begin{align*}
    S^{m-1}=\{\mathbf u\in\R^m:\|\mathbf u\|=1\}.
\end{align*}
These functions are also called hyperspherical harmonics.
They play the same role as the familiar $Y_{\ell q}(\theta,\phi)$: they separate angular dependence into orthogonal modes of increasing degree.
We use $m$ for the dimension and $q$ for a mode label, so $m$ is not the magnetic quantum number of the usual three-dimensional notation.
In our problem, $m=\dim\mathcal V$, the dimension of the teacher span, rather than the input dimension $N$.
We assume $m\geq2$ throughout this appendix.

\subsection{From harmonic polynomials to angular modes}
One way to recognize an ordinary spherical harmonic is to multiply it by the appropriate radial power: $r^\ell Y_{\ell q}$ is a homogeneous harmonic polynomial in three Cartesian coordinates.
The same construction works in any dimension.
Let $Q_\ell$ be a polynomial satisfying
\begin{align*}
    Q_\ell(r\mathbf u)=r^\ell Q_\ell(\mathbf u),
    \qquad \Delta_{\mathbf z}Q_\ell(\mathbf z)=0.
\end{align*}
The first condition says that every monomial has total degree $\ell$; the second says that the polynomial is harmonic.
Its restriction $Y(\mathbf u)=Q_\ell(\mathbf u)$ to $S^{m-1}$ is a generalized spherical harmonic of degree $\ell$.
The polynomial $Q_\ell$ contains both radial and angular dependence, while $Y$ contains only the angular part.

To connect this definition to the familiar angular eigenvalue equation, write $\mathbf z=r\mathbf u$.
The Euclidean Laplacian separates as
\begin{align*}
    \Delta_{\mathbf z}
    =\frac{\partial^2}{\partial r^2}
    +\frac{m-1}{r}\frac{\partial}{\partial r}
    +\frac{1}{r^2}\Delta_{S^{m-1}},
\end{align*}
where $\Delta_{S^{m-1}}$ is the angular Laplacian, or Laplace--Beltrami operator, on the unit sphere.
Applying it to $Q_\ell(r\mathbf u)=r^\ell Y(\mathbf u)$ gives
\begin{align*}
    0=\Delta_{\mathbf z}Q_\ell
    =r^{\ell-2}\left[\ell(\ell+m-2)Y+\Delta_{S^{m-1}}Y\right].
\end{align*}
Thus,
\begin{align}
    -\Delta_{S^{m-1}}Y=\ell(\ell+m-2)Y.
    \label{eq:appendix-harmonic-eigenvalue}
\end{align}
For $m=3$, this is exactly the usual eigenvalue $\ell(\ell+1)$.
For $m=2$, the sphere is a circle, the angular Laplacian is $\partial_\theta^2$, and the positive-degree modes are $\cos(\ell\theta)$ and $\sin(\ell\theta)$.

\subsection{Orthogonality and the number of modes}
As in three dimensions, the degree does not specify a unique harmonic.
The dimension of the degree-$\ell$ harmonic space is
\begin{align*}
    d_{m,\ell}
    =\binom{m+\ell-1}{\ell}-\binom{m+\ell-3}{\ell-2},
\end{align*}
where the second term is taken to be zero for $\ell<2$.
This gives $d_{3,\ell}=2\ell+1$, the usual number of magnetic modes.
On the circle, there is one constant mode and two modes at each positive degree.
We can choose a real orthonormal basis $Y_{\ell q}$, with $q=1,\ldots,d_{m,\ell}$, such that
\begin{align*}
    \int_{S^{m-1}}Y_{\ell q}(\mathbf u)Y_{\ell' q'}(\mathbf u)\,d\Omega(\mathbf u)
    =\delta_{\ell\ell'}\delta_{qq'}.
\end{align*}
Here $d\Omega$ is the usual surface-area measure.
Different degrees are orthogonal because the angular Laplacian is self-adjoint and their eigenvalues in \eqref{eq:appendix-harmonic-eigenvalue} are distinct.
Together, these modes form a complete basis for square-integrable functions on the sphere, just as ordinary spherical harmonics do on $S^2$.

\subsection{Zonal harmonics and Gegenbauer polynomials}
For our overlap function, a useful harmonic is one centered on a unit direction $\mathbf e$ and depending only on $t=\mathbf e^\top\mathbf u$, the cosine of the angle from that direction.
Such a function is called a zonal harmonic; on $S^2$, the familiar example is $P_\ell(\cos\theta)$ with the polar axis along $\mathbf e$.
For a function $Z(t)$, the angular eigenvalue equation becomes
\begin{align*}
    (1-t^2)Z''(t)-(m-1)tZ'(t)+\ell(\ell+m-2)Z(t)=0.
\end{align*}
For $m>2$, its degree-$\ell$ polynomial solution, normalized at $t=1$, is
\begin{align*}
    Z_\ell^{(m)}(t)
    =\frac{C_\ell^{(\nu)}(t)}{C_\ell^{(\nu)}(1)},
    \qquad \nu=\frac{m-2}{2},
\end{align*}
where $C_\ell^{(\nu)}$ is a Gegenbauer polynomial.
Thus Gegenbauer polynomials replace Legendre polynomials when the dimension changes.
At $m=3$, $\nu=1/2$ and $Z_\ell^{(3)}=P_\ell$.
At $m=2$, we instead use $Z_\ell^{(2)}(\cos\theta)=\cos(\ell\theta)$, avoiding the degenerate Gegenbauer normalization at $\nu=0$.
In particular,
\begin{align*}
    Z_1^{(m)}(t)=t,
    \qquad Z_3^{(m)}(t)=\frac{(m+2)t^3-3t}{m-1},
\end{align*}
which are the functions used in equation \eqref{eq:H-zonal-harmonics}.
A zonal harmonic about a chosen axis is a particular linear combination of the degree-$\ell$ basis functions, rather than the entire degree-$\ell$ space.

\subsection{Why a cubic has two harmonic degrees}
Polynomial degree and angular harmonic degree need not coincide.
For example, if $Q_1$ is linear, then $\|\mathbf z\|^2Q_1(\mathbf z)$ is a homogeneous cubic, but on the unit sphere it reduces to the degree-one harmonic $Q_1(\mathbf u)$.
More generally, a homogeneous polynomial of degree $p$ has a unique decomposition
\begin{align*}
    P_p(\mathbf z)
    =\sum_{j=0}^{\lfloor p/2\rfloor}\|\mathbf z\|^{2j}Q_{p-2j}(\mathbf z),
\end{align*}
where each $Q_{p-2j}$ is homogeneous and harmonic of the indicated degree.
Restricting to the sphere removes the radial factors, leaving angular degrees $p,p-2,\ldots$.
For $p=3$, this is precisely $P_3=Q_3+\|\mathbf z\|^2Q_1$, whose explicit construction appears in Section~\ref{sec:cubic-spherical-harmonics}.
This explains why cubic teacher features lead to degrees one and three in $H$, even when there are many features or their directions are correlated.

For further background, see Frye and Efthimiou, \href{https://arxiv.org/abs/1205.3548}{\textit{Spherical Harmonics in $p$ Dimensions}}, especially Chapter~4, and the \href{https://dlmf.nist.gov/18.38}{NIST discussion of zonal spherical harmonics}.

\end{document}
